\slajdid{10-s092}
\begin{frame}[label=10-s092]
\gusttekst\setlength{\abovedisplayskip}{3pt}\setlength{\belowdisplayskip}{3pt}\setlength{\abovedisplayshortskip}{2pt}\setlength{\belowdisplayshortskip}{2pt}\begin{columns}[T,onlytextwidth]\column{.35\textwidth}\centering\input{slike/tikz/s090_ecl_odvojena_napajanja.tex}\column{.63\textwidth}
U dinamičkom režimu parazitna kapacitivnost C na izlazu, podrazumevajući da se kolo nalazi u lancu istih takvih kola, se prazni preko otpornosti $R_{B1}$ sa
\[\tau=CR_{B1}\]
\[v_o(t_0^+)=V_{OH}\ i\ v_o(\infty)=V_{EE}\ i\ v_o(t_0+t_{pHL})=\tfrac{V_{OH}+V_{OL}}2\]
dok se puni preko izlaznih bafera koji rade u aktivnom režimu sa
\[\tau=C\frac{R_C}{\beta_F}\]
\[v_o(t_0^+)=V_{OL}\ i\ v_o(\infty)=V_{CC}\ i\ v_o(t_0+t_{pHL})=\tfrac{V_{OH}+V_{OL}}2\]
Zbog male razlike napona $V_{OH}$ i $V_{OL}$ ovi procesi će biti jako brzi.
\end{columns}
\end{frame}
